\documentclass[11pt,a4paper]{article}

\usepackage[T1]{fontenc}
\usepackage[utf8]{inputenc}
\usepackage{lmodern}
\usepackage{amssymb,amsmath,amsthm,mathtools}
\usepackage[margin=2.5cm]{geometry}
\usepackage{parskip}
\usepackage[hidelinks]{hyperref}
\usepackage{microtype}

% Theorem environments (arXiv / journal style)
\theoremstyle{plain}
\newtheorem{theorem}{Theorem}[section]
\newtheorem{proposition}[theorem]{Proposition}
\newtheorem{lemma}[theorem]{Lemma}
\newtheorem{corollary}[theorem]{Corollary}

\theoremstyle{definition}
\newtheorem{definition}[theorem]{Definition}
\newtheorem{example}[theorem]{Example}
\newtheorem{remark}[theorem]{Remark}

% Notation
\newcommand{\Ext}{\mathrm{Ext}}
\newcommand{\Shape}{\mathrm{Shape}}
\newcommand{\Pos}{\mathrm{Pos}}
\newcommand{\Type}{\mathbf{Type}}
\newcommand{\id}{\mathrm{id}}
\newcommand{\Cont}{\mathbf{Cont}}
\newcommand{\rootp}{\mathrm{o}}

\title{Directed Containers as Comonoids\\
  in the Sequential Monoidal Structure on Containers}
\author{Extracted from a machine-checked development}
\date{}

\begin{document}
\maketitle

\begin{abstract}
We define comonoids in the monoidal category of containers under sequential
composition, and show that every directed container determines such a
comonoid. Counit and comultiplication are container morphisms
$\varepsilon : C \Rightarrow I$ and $\delta : C \Rightarrow C \mathbin{\triangleleft} C$,
and the comonoid laws are equations of container morphisms built from the
sequential action on morphisms and the monoidal unitors and associator.
Unfolding these laws recovers the directed-container axioms D1--D5. The proofs
of all stated results have been machine-checked; we record only the
mathematical content.
\end{abstract}

\tableofcontents
\bigskip

%======================================================================
\section{Introduction}
%======================================================================

The category $\Cont$ of containers and container morphisms carries a monoidal
structure $(\Cont,\, \triangleleft,\, I)$ whose tensor is sequential
composition: the extension of $C \mathbin{\triangleleft} D$ is the composite of
the extensions of~$C$ and~$D$, and~$I$ is the unit container (a single shape
with a single position). A \emph{comonoid} in this monoidal category is a
container~$C$ together with morphisms
\[
  \varepsilon : C \Rightarrow I,
  \qquad
  \delta : C \Rightarrow C \mathbin{\triangleleft} C
\]
satisfying the two counit laws and coassociativity, all stated internally to
$\Cont$.

This note shows that every directed container---a container with root,
sub-shape, and shift satisfying D1--D5---gives rise to such a comonoid. The
point of the internal formulation is that the laws never mention extensions or
endofunctors on types: they are equations between container morphisms. Upon
unfolding, one recovers the directed-container axioms on the nose, with the
same partition that appears at the comonad level:
\begin{itemize}
  \item right counit is D3;
  \item left counit is D1 (shapes) and D2 (positions);
  \item coassociativity is D4 (shapes) and D5 (positions).
\end{itemize}
Thus the correspondence between directed containers and comonads on extensions
is an artifact of the monoidal structure on $\Cont$, not merely of passage
through the functor category.

We work in a single universe of types. Proofs are indicated by the phrase
``by lean.''

\paragraph{What this note claims.}
Directed containers are $\triangleleft$-comonoids. Explicitly: the root
supplies a counit $C \Rightarrow I$; sub and shift supply a comultiplication
$C \Rightarrow C \mathbin{\triangleleft} C$; and D1--D5 are precisely the three
comonoid laws for sequential composition. (The converse---that every comonoid
arises this way---is not treated here.)

%======================================================================
\section{Background}
%======================================================================

\begin{definition}[Container morphism]
A \textbf{morphism} $\varphi : C \Rightarrow D$ of containers consists of a
shape map $u : \Shape(C) \to \Shape(D)$ and, for each shape~$s$, a
contravariant position map
$g_s : \Pos(D)(u(s)) \to \Pos(C)(s)$.
\end{definition}

\begin{definition}[Sequential composition]\label{def:seq}
The \textbf{sequential composite} $E \mathbin{\triangleleft} C$ of containers
has shapes $\sum_{t \,:\, \Shape(E)} \bigl(\Pos(E)(t) \to \Shape(C)\bigr)$
and positions over $\langle t,\, f\rangle$ given by
$\sum_{u \,:\, \Pos(E)(t)} \Pos(C)(f(u))$. The unit container~$I$ has a
single shape and a single position.
\end{definition}

We assume the monoidal structure on $(\Cont,\, \triangleleft,\, I)$:
associator $\alpha$, left and right unitors $\lambda$ and $\rho$, and the
action of $\triangleleft$ on morphisms (denoted $\triangleleft_2$ or
$\mathrm{seq}_2$ when we need the binary map on arrows). Composition of
container morphisms is written in diagrammatic order when convenient, so that
$\delta ; (\varepsilon \mathbin{\triangleleft} C)$ means
$(\varepsilon \mathbin{\triangleleft} C) \circ \delta$ in applicative order.

\begin{definition}[Directed container]
A \textbf{directed container} is a container equipped with
\begin{align*}
  \rootp &\;:\; (s : S) \to P(s), \\
  \downarrow &\;:\; (s : S) \to P(s) \to S, \\
  \oplus &\;:\; (s : S) \to (p : P(s)) \to P(s \downarrow p) \to P(s),
\end{align*}
satisfying axioms D1--D5 (sub under root; shift by root; shift by sub-root;
nesting of subs; associativity of shift).
\end{definition}

%======================================================================
\section{Technical lemmas on morphisms and transport}
%======================================================================

Because some directed-container axioms are only propositional equalities of
shapes, comparing container morphisms whose shape maps agree only
propositionally requires pointwise extensionality and a computation of
position transport under sequential composition.

\begin{lemma}[Pointwise extensionality of morphisms]\label{lem:ext-eq}
Let $f,g : \Shape(C) \to \Shape(D)$ and let $f_p$, $g_p$ be position maps
over~$f$ and~$g$ respectively. If $h_s : f(s) = g(s)$ for every~$s$ and
$f_p(s,p) = g_p\bigl(s,\, h_s \mathbin{\triangleright} p\bigr)$ for every
position~$p$, then the morphisms $\langle f,\, f_p\rangle$ and
$\langle g,\, g_p\rangle$ are equal.
\end{lemma}

\begin{proof}
By lean.
\end{proof}

\begin{lemma}[Cancelation of cast round-trips]\label{lem:cast-symm}
For a family $P$ and an equality $h : a = b$, transporting a value of
$P(a)$ along~$h$ and back along $h^{-1}$ returns the original value.
\end{lemma}

\begin{proof}
By lean.
\end{proof}

\begin{lemma}[Sequential position transport]\label{lem:seq-pos-transport}
Let $E$ and $C$ be containers, let $t$ be a shape of~$E$, and let
$f,g : \Pos(E)(t) \to \Shape(C)$ be inner-shape assignments with
$h_u : f(u) = g(u)$ for every outer position~$u$. Transport of a position
$\langle u,\, q\rangle$ of $E \mathbin{\triangleleft} C$ over
$\langle t,\, f\rangle$ along the induced shape equality yields
\begin{equation}
  \bigl\langle u,\;
    (h_u \mathbin{\triangleright} q) \bigr\rangle
\end{equation}
over $\langle t,\, g\rangle$: the outer position is fixed and only the inner
position is transported.
\end{lemma}

\begin{proof}
By lean.
\end{proof}

%======================================================================
\section{Comonoids in $(\Cont, \triangleleft, I)$}
%======================================================================

\begin{definition}[Comonoid]\label{def:comonoid}
A \textbf{comonoid} in $(\Cont,\, \triangleleft,\, I)$ consists of a
container~$C$ together with morphisms
\begin{equation}
  \varepsilon : C \Rightarrow I,
  \qquad
  \delta : C \Rightarrow C \mathbin{\triangleleft} C,
\end{equation}
such that the following equations of morphisms hold:
\begin{align}
  \delta ; \bigl(\varepsilon \mathbin{\triangleleft} C\bigr)
  &\;=\;
  \lambda_C^{-1}
  &&\text{(left counit),} \\[4pt]
  \delta ; \bigl(C \mathbin{\triangleleft} \varepsilon\bigr)
  &\;=\;
  \rho_C^{-1}
  &&\text{(right counit),} \\[4pt]
  \delta ; \bigl(\delta \mathbin{\triangleleft} C\bigr) ; \alpha_{C,C,C}
  &\;=\;
  \delta ; \bigl(C \mathbin{\triangleleft} \delta\bigr)
  &&\text{(coassociativity).}
\end{align}
\end{definition}

\begin{remark}
These are the standard comonoid axioms in a monoidal category, written with
the sequential tensor and its coherence isomorphisms.
\end{remark}

%======================================================================
\section{From directed containers to comonoids}
%======================================================================

\begin{definition}[Counit from a directed container]\label{def:counit}
Let $D$ be a directed container. The \textbf{counit}
$\varepsilon : D \Rightarrow I$ is the container morphism with unique shape
map and position map
\begin{equation}
  \varepsilon_p(s,\, \ast) \;=\; \rootp(s).
\end{equation}
\end{definition}

\begin{definition}[Comultiplication from a directed container]\label{def:comult}
The \textbf{comultiplication}
$\delta : D \Rightarrow D \mathbin{\triangleleft} D$ is the container morphism
with
\begin{align}
  \delta_S(s) &\;=\; \bigl\langle s,\; \lambda p.\, s \downarrow p \bigr\rangle,
  \\[4pt]
  \delta_p\bigl(s,\, \langle p,\, q\rangle\bigr)
  &\;=\; p \oplus q.
\end{align}
\end{definition}

\begin{theorem}[Directed containers are comonoids]\label{thm:to-comonoid}
For every directed container~$D$, the morphisms $\varepsilon$ and $\delta$ of
Definitions~\ref{def:counit} and~\ref{def:comult} satisfy the three comonoid
laws of Definition~\ref{def:comonoid}. Thus $D$ is a
$\triangleleft$-comonoid.
\end{theorem}

\begin{proof}
By lean.
\end{proof}

\begin{remark}[How the axioms match]
\begin{enumerate}
  \item \textbf{Right counit.} Both sides have definitionally equal shape
        maps; the position equality is exactly D3.
  \item \textbf{Left counit.} Shape equality is D1; after sequential position
        transport (Lemma~\ref{lem:seq-pos-transport}), the position equality
        is D2.
  \item \textbf{Coassociativity.} Shape equality is D4 (pointwise in outer
        and inner positions); position equality is D5 after cancelling a cast
        round-trip (Lemma~\ref{lem:cast-symm}) and applying sequential
        transport twice.
\end{enumerate}
\end{remark}

%======================================================================
\section{Discussion}
%======================================================================

\paragraph{What has been established.}
The directed-container structure is precisely the data of a comonoid in
$(\Cont,\, \triangleleft,\, I)$, at least in the forward direction. The
comonoid laws are equations of container morphisms; their unfolding is the
classical list D1--D5. In particular, one need not pass through
$[\mathbf{Set},\, \mathbf{Set}]$ to state or verify the comonoid equations.

\paragraph{What is not established here.}
The converse---that every $\triangleleft$-comonoid arises from a directed
container---is not formalised in the source development. One expects it by
the same representation and shape-collapse arguments that recover directed
data from a container comonad on the extension, but that direction is left
open.

\paragraph{Machine checking.}
Every definition and theorem statement above corresponds to a declaration in a
machine-checked formalisation. The phrase ``by lean'' indicates that the proof
has been verified by a proof assistant.

%======================================================================
\section*{Acknowledgements}
%======================================================================

The mathematical content follows the monoidal structure of sequential
composition of containers and the directed-container theory of Ahman, Chapman,
and Uustalu. The presentation here is a prose rendering of a formal
development intended for human review of the theorem statements.

\begin{thebibliography}{9}

\bibitem{ahman2014}
D.~Ahman, J.~Chapman, and T.~Uustalu.
\newblock When is a container a comonad?
\newblock \emph{Logical Methods in Computer Science}, 10(3), 2014.

\bibitem{abbott2005}
M.~Abbott, T.~Altenkirch, and N.~Ghani.
\newblock Containers: constructing strictly positive types.
\newblock \emph{Theoretical Computer Science}, 342(1):3--27, 2005.

\end{thebibliography}

\end{document}
